\documentclass[11pt,a4paper]{article}
\usepackage{iftex}
\ifXeTeX
  \usepackage{fontspec}
  \setmainfont{texgyretermes}[Extension=.otf, UprightFont=*-regular, BoldFont=*-bold,
    ItalicFont=*-italic, BoldItalicFont=*-bolditalic]
\else
  \usepackage[T1]{fontenc}
  \usepackage[utf8]{inputenc}
  \usepackage{newtxtext}
\fi
\usepackage[turkish]{babel}
\usepackage[margin=2cm]{geometry}
\usepackage{booktabs,tabularx,xltabular,array,enumitem,xcolor}
\usepackage[hidelinks]{hyperref}
\usepackage{xurl}
\urlstyle{tt}
\newcolumntype{Y}{>{\raggedright\arraybackslash}X}
\setlist{nosep,leftmargin=*}
\setlength{\parskip}{0.4em}
\setlength{\parindent}{0pt}

\title{Ek A -- GenAI Kullanım Beyanı\\[2pt]\large Ara Rapor: ALL/AML Lösemi Alt Tip Sınıflandırması}
\author{Enes AY -- 503020260092\\\small Eskişehir Osmangazi Üniversitesi, Fen Bilimleri Enstitüsü, Yapay Zeka Programı}
\date{Kasım 2026}

\begin{document}
\maketitle

\section*{1. Beyan}
Bu projenin ara rapor aşamasında Anthropic'in \textbf{Claude Opus 5.5} modeli (Claude Code masaüstü
uygulaması) yardımcı araç olarak kullanılmıştır. Projenin kapsamı, bonus hedefleri, veri setleri,
rapor dili ve yöntemsel kararlar tarafımca belirlenmiş; model bu kararları uygulamada, anlamadığım
konuları açıklamada, sonuçları test etmede ve konuda derinleşmede kullanılmıştır. Kod Claude Code ile
yazılmış, tarafımca çalıştırılarak doğrulanmıştır. Literatür taraması ve makalelerin incelenmesi
Claude Opus 5.5 ile birlikte yapılmış, temel makaleler tarafımca okunmuştur. Modelin ürettiği tüm
içerik tarafımca gözden geçirilmiştir; raporun ve kodun içeriğine ilişkin tüm sorumluluk bana aittir.

\section*{2. Kullanım tablosu}
{\small
\begin{xltabular}{\textwidth}{@{}>{\raggedright\arraybackslash}p{2.2cm} Y >{\raggedright\arraybackslash}p{2.9cm} Y Y@{}}
\toprule
Araç / sürüm & Amaç & Bölüm / dosya & Öğrencinin katkısı & Nasıl doğrulandı \\
\midrule
\endhead
Claude Opus 5.5 & Rehber isterlerinin maddeleştirilmesi ve proje planı & tüm proje &
Kapsamı, bonus hedeflerini ve takvimi belirledi; ara raporun Türkçe, final makalesinin İngilizce olmasına karar verdi &
İsterler rapor bölümleriyle tek tek eşlendi (\url{TESLIM_KONTROL.md}) \\
Claude Opus 5.5 (Claude Code) & Veri indirme, yükleme, ön işleme ve EDA kodu & \url{scripts/}, \url{src/leukemia/} &
Gereksinimleri belirledi; daha güçlü ve bağımsız veri setleri aranmasını istedi; betikleri kendi bilgisayarında çalıştırdı &
EDA betiği yeniden çalıştırıldı, sayılar raporla aynı çıktı; kodda etiket ve sınıf sayısı kontrolleri (\texttt{assert}) \\
Claude Opus 5.5 & Literatür taraması ve makalelerin incelenmesi & \url{reports/literature/}, rapor Bölüm II &
Son beş yılın kaynaklarına ağırlık verilmesini ve iddiaların kaynak metinle denetlenmesini istedi; temel makaleleri okudu &
DOI'ler Crossref ile doğrulandı; rapordaki iddialar kaynak metinlerle karşılaştırıldı (\url{claim_audit.md}) \\
Claude Opus 5.5 & Kavramların açıklanması ve konuda derinleşme & -- &
Anlamadığı kavramları (sızıntı, iç içe CV, AUC ve karar eşiği vb.) sordu; katkıların bir bildiri için yeterliliğini sorguladı &
Kavramlar videoda kendi sözleriyle anlatıldı \\
Claude Opus 5.5 & Fizibilite pilotu (planlanan ölçümlerin denenmesi) & \url{scripts/run_feasibility.py}, rapor Bölüm V &
Raporun final ve makale için yeterliliği göstermesini istedi &
Sabit ayarlar; kilitli test setlerine dokunulmadı \\
Claude Opus 5.5 & Rapor ve sunum metninin yazımı ve Türkçeye çevirisi & \url{reports/ara_rapor/}, \url{reports/sunum/} &
Rapor dilini ve yapısını belirledi; düzenleme için Word kopyası istedi &
Rapordaki sayılar betik çıktılarından otomatik geliyor \\
Claude Opus 5.5 & Video için anlatım taslağı & \url{reports/video/} &
Videoyu kendi sesiyle ve kendi cümleleriyle kaydetti & Video süresi ve bağlantı erişimi kontrol edildi \\
\bottomrule
\end{xltabular}}

\section*{3. Önemli prompt'ların özeti}
\begin{itemize}
\item Ders rehberindeki tüm şartları (bonuslar dahil) karşılayan bir proje planı ve uygulaması istendi.
\item Kaynakların gerçekten okunup okunmadığı soruldu; bunun üzerine iddialar kaynak metinlerle denetlendi.
\item Raporun final ve uluslararası makale için yeterliliği, daha güçlü veri setleri ve güncel kaynaklar istendi.
\item Projenin amacı, çalışma konusunun nasıl belirlendiği ve kullanılan terimler soruldu.
\item Katkıların bir uluslararası bildiri için ne kadar güçlü olduğu soruldu.
\end{itemize}

Doğrulama sırasında modelin hataları yakalanıp düzeltilmiştir: yanlış bir DOI, yanlış tanımlanan bir
platform (GSE7186), MILE verisinin ölçeği ve gereğinden geniş bir özgünlük iddiası.

\vspace{1em}
\noindent Enes AY (503020260092) \hfill Kasım 2026
\end{document}
